\section{Scalar expansions and the tree inequality}
\label{ell:scalar-calculus}

\begin{lemma}[Finite-part integration]
\label{ell:finite-parts}
Integration of the coefficient fields constructed in
Lemma~\ref{ell:one-length} preserves exponential-polynomial
length expansions for $F$ and $J$.  Each resulting coefficient
is a finite sum of core or root-tail integrals, finite-part
integrals on opened ends, and evaluations of finite tail
coefficients, on the corresponding partial disconnection.
These operations are analytic and continuous in sufficiently
high finite tail precision.
\end{lemma}

\begin{proof}
For clarity fix a normalization of finite parts.  Let an
integrand density on a complete opened end, already integrated
over its cross-section, be $f(\tau)d\tau$.  Subtract a fixed
cutoff times every nondecaying term of its finite tail,
and denote this finite sum by
$T(\tau)=\sum_{\alpha\ge0,j}c_{\alpha j}e^{\alpha\tau}\tau^j$.
Let $Q_{\alpha j}$ be a chosen polynomial-exponential primitive
of $e^{\alpha\tau}\tau^j$, with zero additive constant.
The remainder $f-\zeta T$ is integrable once its tail rate is
positive.  There is then a uniquely normalized constant $C_f$
such that
\begin{equation}
 \int_{\tau_0}^{L}f(\tau)\,d\tau
 =C_f+\sum_{\alpha\ge0,j}c_{\alpha j}Q_{\alpha j}(L)
             -\int_L^\infty(f-T)(\tau)\,d\tau
 \label{ell:finite-part-formula}
\end{equation}
when $L$ is beyond the cutoff.  Define $C_f$ by integrating
$f-\zeta T$ and the fixed compact correction.  The final
integral uses the exact tail after subtraction and can also
be expanded by subtracting its finitely many decaying terms.
If the remaining rate is $\eta>0$, its integral is bounded
by $C_\eta$ times that finite tail norm, by Cauchy--Schwarz
against $e^{-\eta u}$.  This proves continuity.  It also proves
analyticity by locally uniform dominated integration; the
coefficients themselves are analytic.  On complex contours,
use $d\tau$ and the contour-independence argument from
Lemma~\ref{ell:analytic-family}.  A strict exponential margin
pays the polynomial length of the contour.

Apply this formula to each Taylor coefficient of the densities
in Equation~\eqref{ell:observables}.  They are finite analytic
differential operations on the base field, or finite Taylor
operations on appended fields.  The interior Ricci-square
density is scale neutral.  A scalar-volume density has the
additional factor $a_v^2e^{\pm2\tau}$.  On a newly opened
child half its endpoint factor is accompanied by $e^{-4L}$.
Thus all primitives still have integer exponential powers
and polynomial logarithmic factors, with the stated grid
after multiplication by the factors in
Equation~\eqref{ell:length-terms}.

Core integration is continuous in the basic Sobolev norm.
Root-tail integration is continuous in the fixed root weights,
using $q_0>1$ for the subtracted density and $q_1<q_0$ for
the weighted Ricci square.  On other finite edges the child
scale factors still accompany the $e^{2\tau}$ radial factor,
so the boundedness estimates following
Equation~\eqref{ell:observables} apply there.  Their integrands
have the requisite decaying field factors.  Evaluation of a
finite part or tail coefficient on an older opened end has
no loss in the new length: it is continuous in a fixed
sufficient finite tail norm, where
Lemma~\ref{ell:one-length} retained that new-length factor.

To reach a scalar cutoff $B$, choose a tag cutoff strictly
greater than $B+\nu+2+\delta$ for some $\delta>0$, then choose
the tail precision higher still if required by the finitely
many auxiliary growth exponents.  The allowance $\nu$ pays
the basic weight, $2$ pays the worst radial factor, and the
strict surplus pays every polynomial factor.  The actual
child factor $e^{-4L}$ only improves this bound.  Thus the
formal integration and endpoint truncation have genuine
exponentially better remainders.
\end{proof}

\begin{proposition}[All-order scalar expansion]
\label{ell:all-orders}
The functions $F,J$ satisfy the analytic expansion hypotheses
of Theorem~\ref{thm:path-selection}.  Their length exponents
belong to the discrete nonnegative grid $\lambda\N\cup\{0\}$.
In any length, and recursively after taking any finite list
of coefficients, expansions have exponentially better
remainders at every cutoff.  All recursively obtained
coefficients are analytic on a common fixed-width complex
tail band of the remaining lengths and one common open
complex neighborhood of the ordinary parameters.
The neighborhood does not decrease with the expansion order.
\end{proposition}

\begin{proof}
Lemmas~\ref{ell:one-length} and~\ref{ell:finite-parts} give an
expansion in one length. Its coefficients are finite constructions
of the same kind: triangular field equations with finite free
asymptotic data, followed by core and root integrals, finite parts,
and tail-coefficient evaluations. On every still-finite seam the
added fields have homogeneous matching and zero critical tests.
Opening the next edge therefore uses the same lemmas; variations
of conditions at older opened ends expand their linear free-coefficient
equations in triangular order.

We check that each such coefficient is bounded uniformly in the
lengths still finite. Fix its entire finite construction and a
sufficient finite tail precision. A prescribed growing tail is cut
off on an already opened end in
Equation~\eqref{ell:finite-growth-inverse}. The correction is obtained
from the fixed basic inverse, so it still decays on every finite
half-end. Repeating the finitely many triangular solves preserves
uniform field and tail bounds. Polynomial prescriptions in earlier
tail coefficients are bounded by the preceding steps, and the two
decaying factors on still-finite halves pay their strengthened source
norms. Constants may depend on this coefficient, its growth and its
chosen precision, but not on the remaining lengths or parameters
in a fixed smaller polydisc.

The scalar operations preserve these bounds. Core and root integrals
are continuous in the fixed norms. Tail-coefficient evaluation and
finite-part integration on an older opened end are bounded maps at
the chosen sufficient precision, by Lemma~\ref{ell:finite-parts}.
On a still-finite edge keep the physical scale factors belonging to
unexpanded lengths. Scale-neutral Ricci-square terms have two
decaying factors and hence a uniformly bounded cylindrical integral.
For scalar-volume terms, the parent punctured half has the integrable
factor $a_p^2e^{-2\tau}$, while the child half satisfies
\[
 a_c^2\int_{\tau_0}^{L_e}e^{2\tau}\,d\tau
 \le C a_p^2e^{-2L_e},\qquad a_c=a_pe^{-2L_e}.
\]
The bounded normalized fields multiply these factors. On complex
contours the corresponding bounds use real parts; polynomial contour
costs are absorbed by decay on the bent portion or by the displayed
child factor. Thus each recursively extracted scalar coefficient
has the required uniform bound on quadratic regions, with thresholds
and constants allowed to depend on its finite construction.

Every scalar exponent lies on $\lambda\Z$. Group equal exponents.
For fixed real remaining parameters, boundedness as the expanded
length tends to infinity excludes a nonzero negative exponent and
a positive polynomial degree at exponent zero: the first such term
would dominate all later terms. The recursive bounds just proved
permit the same argument after any previous coefficient extraction.
Finite expansions are unique, since a nonzero polynomial cannot
decay exponentially after multiplication by the exponential of the
smallest differing exponent. Applying this at real remaining
parameters and then using analytic continuation makes the formulas
from different cutoffs agree. This gives consistent expansions on
the nonnegative grid.

We give the common-domain argument separately, since merely
having an expansion to every finite order would not imply it.
For each of the finitely many partial disconnections, fix its
small solution and its augmented basic inverse from
Lemma~\ref{ell:analytic-family}.  Take once the minimum of
their parameter radii and the maximum of their length thresholds,
and restrict to
a smaller polydisc compactly contained in all of them.
Their common quadratic regions include a band
\begin{equation}
 \Re L_e>A_*,\quad |\Im L_e|<c_*;
 \qquad |z-z_*|<\rho_*.
 \label{ell:common-domain}
\end{equation}
Every coefficient of every finite order is constructed on
this same domain by finite polynomial primitives, cutoff
tails, and the same basic inverse, as in
Equation~\eqref{ell:finite-growth-inverse}.  A large growing
coefficient never becomes the input of a new nonlinear
contraction.  Its correction is an unrestricted-size linear
solve.  Increasing weights in Lemma~\ref{ell:flat-tail} uses
only constant root kernels to identify tails; it does not
introduce a new parameter-dependent global inverse.  Tail
limits are locally uniform on this domain, and finite parts
are locally uniformly integrable after the finite subtraction.
These observations prove analyticity on
Equation~\eqref{ell:common-domain} for every degree.

For the remainder bounds on quadratic regions it is permissible
to increase length thresholds and constants with the finite
precision and finite auxiliary system.  The ordinary-parameter
polydisc remains the one fixed above.  Choose a fixed slightly
larger polydisc still inside the basic inverse domain whenever
Cauchy estimates are needed; this choice is also made once.
No intersection of infinitely many decreasing parameter
neighborhoods occurs.  On complete opened ends used for
coefficient formulas one may use real contours to establish
this common-band analyticity, then the contour estimates to
obtain the stronger asymptotic bounds.

Finally consider a subset of lengths with
$\Re L_i=c_i U+o(U)$, $c_i>0$.  There are only finitely many
grid multi-indices with $\sum_i d_ic_i\le B$ for any fixed
$B$.  Expand the first length sufficiently far that its
remainder has order greater than $(B+\epsilon)/c_1$.
For each of its finitely many retained coefficients, expand
the next length sufficiently far to put the resulting
weighted remainder above $B+\epsilon$, and continue.
Uniform coefficient bounds in the still unexpanded lengths
make this a legitimate finite procedure.  Choose the
strict surplus before making the finitely many choices.
The $o(U)$ changes of the length ratios and all logarithmic
polynomials consume less than that surplus for large $U$.
This gives the required grouped weighted expansions with
strictly better remainders, including after previous
coefficients have been taken.  These are precisely the
remaining expansion conditions in
Theorem~\ref{thm:path-selection}.
\end{proof}

\subsection{The differential estimate and the contradiction}
\label{ell:conclusion}

\begin{lemma}[Differential of the scalar observable]
\label{ell:differential}
On the real small family, for some $c>0$,
\begin{equation}
 |dF|\le C\sqrt J\left(\sum_\alpha|dz_\alpha|
                 +\sum_e e^{-cL_e}|dL_e|\right).
 \label{ell:dF}
\end{equation}
\end{lemma}

\begin{proof}
For ordinary-parameter derivatives the uniform linearized
inverse gives the fixed weighted field bounds.  Pair the
first variation with the first Ricci-square term of $J$ on
the physical interior.  Its dual metric-variation $L^2$
norm is uniformly bounded: core volumes scale by $a_v^4$,
and the half-end volume calculation uses
$a_v^4e^{\pm4\tau}d\tau$, with the factor $e^{-8L_e}$ on
a child side.  The scalar part of the Einstein tensor costs
only a dimension-dependent multiple of $|\Ric|$.
On the root use instead the weighted part of $J$ and
Cauchy--Schwarz.  The dual integral is bounded by
\[
 C\int_{r_0}^{\infty}r^{-2q_1}r^{-2q_0}r^3\,dr<\infty,
\]
since $q_0+q_1>2$.  This proves the $dz$ part.

Now differentiate one length $L_e$, initially at fixed vertex
and end coordinates.  The normalized bulk equations and
the compact tests have no dependence on this length.  Only
the seam data acquire an inhomogeneity: derivatives of the
prescribed $e^{-\lambda L_e}z_e$, and endpoint normal
derivatives of the solution in its matching conditions.
The solution has the stronger joining weight $\nu'>\nu$.
The trace theorem, using one fewer derivative if necessary,
therefore bounds these differentiated data in the basic
weighted norms by
\[
 C\bigl(e^{-(\nu'-\nu)L_e}
             +e^{-(\lambda-\nu)L_e}\bigr)|dL_e|.
\]
The linearized inverse gives this same bound throughout
the normalized field.  Its first-variation pairing is
bounded by $C\sqrt J e^{-c_1L_e}|dL_e|$ for any fixed
$0<c_1<\nu'-\nu$, by the preceding dual-norm estimates.

To interpret this derivative on a fixed smooth manifold,
move the endpoints in a collar of fixed logarithmic width
at the seam.  The additional transport variation is bounded
in relative metric norm and supported at physical radius
$\rho\le Ca_pe^{-L_e}$.  Its $L^2$ norm is therefore at most
$Ca_p^2e^{-2L_e}|dL_e|$.  Moreover the physical squared scale
of every descendant changes at a bounded relative rate,
since $\partial_{L_e}a_v^2=-4a_v^2$ there.  The whole
descendant region has volume at most
$Ca_p^4e^{-4L_e}$.  To verify this last assertion, the first
child exterior half has volume
\[
 Ca_c^4\int^{L_e}e^{4\tau}\,d\tau
       \le Ca_p^4e^{-4L_e};
\]
its core and each further subtree are smaller by their
additional scale factors, and the tree is finite.  The
descendant scaling variation thus has the same exponentially
small $L^2$ bound.  Pair both variations with the unweighted
Ricci-square term of $J$.  Choosing
$0<c<\min(\nu'-\nu,2)$ proves
Equation~\eqref{ell:dF}.  All identifications here use real
lengths; no positivity of a complex bilinear square is used.
\end{proof}

\begin{proof}[Proof of Theorem~\ref{thm:tree-inequality}]
The case $M=0$ is Lemma~\ref{ell:functional}.  Suppose the
assertion fails on a subsequence with $M>0$.  Then
$|E(g)|\ge\epsilon M$ for some fixed $\epsilon>0$.
By Lemma~\ref{ell:comparison} and
Equation~\eqref{ell:J-small}, late terms satisfy
\[
 F\ne0,\qquad F^2\ge c_2J,\qquad F\longrightarrow0,
\]
for some $c_2>0$.  The finite parameter vectors have an
interior subsequential limit: all the earlier smallness
choices can be strict inside the parameter polydisc.
On real parameters, $J\ge0$ by Equation~\eqref{ell:observables}.
Proposition~\ref{ell:all-orders} supplies admissibility, and
Lemma~\ref{ell:differential} supplies Equation~\eqref{ell:dF}.
The resulting sequence therefore satisfies every hypothesis
excluded by Corollary~\ref{path:variation}, a contradiction.
This proves Equation~\eqref{ell:little-oh}.
\end{proof}
